\section{Introduction}
\label{sec:introduction}

Recent advances in text-to-speech, voice cloning, and speech generation have made audio deepfakes increasingly realistic and accessible, turning audio deepfake detection into an important problem in multimedia security. In response, the community has developed a broad range of benchmarks and datasets that substantially improve evaluation across generators, tasks, languages, and domains~\cite{liu2023asvspoof,yi2024add,frank2021wavefake,khalid2021fakeavceleb,muller2022generalize,zhang202partial,sun2023ai,muller2024mlaad,ma2024cfad,wang2026asvspoof,xie2025codecfake,bhagtani2025diffssd,ciobanu2026xmad,wang2025audeter}. These resources have considerably strengthened evaluation on the generation side of audio deepfakes.

However, realistic misuse rarely ends at synthesis. In practice, a forged clip may be uploaded to an online platform, transmitted through a telephony channel, or played aloud and re-recorded before it reaches a listener or an automated system. The audio encountered at deployment is therefore often not pristine generated speech, but speech after post-generation delivery. This creates a mismatch between common benchmark conditions and realistic use scenarios. Some recent studies have started to consider communication effects, codec distortions, or replay-like conditions~\cite{shi2025benchmarking,muller2025replay,delgado2026presentation}, but they mostly examine delivery through isolated effects. While useful, such settings do not fully reflect the fact that post-generation delivery is often a structured process in which multiple operations are combined, parameterized, and ordered. Platform re-encoding, telephony transmission, and replay-like recapture may affect forensic cues differently, and their impact may depend on the particular delivery route rather than on any single operation alone.

These considerations suggest the need for a benchmark that treats post-generation delivery as a structured evaluation dimension: one that spans multiple realistic delivery families, represents delivery as ordered chains rather than isolated perturbations, and supports controlled analysis under matched conditions. Such control is important because, without matched clean parents, delivery effects can be entangled with differences in linguistic content, speaker identity, or source audio quality, making it difficult to attribute detector behavior specifically to delivery.

To this end, we introduce ChainBench-ADD, a delivery-aware dataset and benchmark for audio deepfake detection. It models post-generation delivery through reusable operators, ordered templates, and realized chains across five delivery families: direct, platform-like, telephony, simulated replay, and hybrid. Each delivered sample remains linked to a clean bona fide or spoof parent under matched-parent control, enabling controlled comparison across delivery scenarios. Figure~\ref{fig:teaser} summarizes the benchmark design.

The current release contains over 941k waveforms derived from 55,813 parents. From the shared metadata, we define five evaluation tasks: in-chain detection, three matched local interventions (operator substitution, parameter perturbation, and order swap), and lineage-based delivery robustness. A leave-one-template-out protocol further tests transfer to unseen templates within a family. Together, these components support more realistic and controlled delivery-aware evaluation for audio deepfake detection.

The main contributions of this work are as follows:
\begin{itemize}
    \item We introduce ChainBench-ADD, a large-scale delivery-aware dataset and benchmark for audio deepfake detection that models post-generation delivery as structured chains across five realistic families through reusable operators, ordered templates, and realized chains.
    \item We design a matched-parent, metadata-driven benchmark framework that isolates delivery effects under fixed transcript and speaker conditions and supports controlled, protocol-aware evaluation within a unified metadata schema.
    \item We release a protocol-aware resource with explicit lineage annotations, benchmark statistics, and consistency validation, and show through baseline studies that it exposes family-level, local, and path-wise robustness patterns that conventional pooled evaluation cannot capture.
\end{itemize}